%%%%%%%%%%%%%%%%%%%%%%%%%%%%%%%%%%%%%%%%%%%%%%%%%%%%%%%%%%%%%%%%%%%%%%
% How to Write a Scientific Thesis
% Author: Dr.-Ing. Mark Schutera
% Based on Tufte-style layout
\makeindex
%%%%%%%%%%%%%%%%%%%%%%%%%%%%%%%%%%%%%%%%%%%%%%%%%%%%%%%%%%%%%%%%%%%%%%

\documentclass{../sharedAssets/tufte}

%%
% Book metadata
\title{Thesis Road Map}
\author[\defaultauthor]{\defaultauthor}
\publisher[\defaultpublisher]{\defaultpublisher}

% Load optional fonts
\loadoptionalfonts

% Additional packages
\usepackage{pifont}
\usepackage{amsmath}
\usepackage{soul}
\usepackage{textcase}
\usepackage{enumitem}
\usepackage{tikz}
\usepackage{pgfplots}
\usepackage{float}
\usepackage{algorithm}
\usepackage{algpseudocode}

\begin{document}



\makestandardfrontmatter
% Dedication
% Alternative quotes:
\makededication{\openepigraph{“Let others follow, if they can! Let them a journey new begin.”}{Bilbo Baggins}}

% \tableofcontents




% Introduction
\makeintroduction{%
  \newthought{You have very little time} to write your thesis and complete the required development and research.
  Alongside that, there is bureaucracy and paperwork to handle.

  \newthought{This roadmap} aims to simplify the process by providing clear steps, points of contact, and the key information you need.
  It outlines the thesis process from eligibility to defense.
  Think of this as a hub, with spokes leading to more detailed resources.

  
  \newthought{Feedback is invaluable} in this process. Names, links, procedures, and requirements may change over time.
  If you have suggestions, corrections, or additional tips, please contribute.
  This helps keep the roadmap accurate and useful for everyone. As many did before you, pay it forward.
  }


\vfill
\par\smallcaps{Contribute - } This project is open for contributions, 
via GitHub \url{https://github.com/Quillstacks/ThesisWork/ThesisRoadmap} to report issues, suggest improvements, and submit pull requests.







%% Start the main matter
\mainmatter
\chapter*{Roadmap for Bachelor's Thesis}


\bigskip
\section{Overview of the Artifacts}
\newthought{Three artifacts scaffold the 12-week writing period} from the beginning of May until the end of July. Each artifact is a milestone that structures your progress and feeds directly into the final thesis.
\marginnote{\textbf{Timeline:} May 1 -- July 31.}

\begin{enumerate}
    \item \textbf{Research Proposal} (prior to final topic registration) -- Literature research and problem statement. Introduction, motivation, gap in the state of the art, research questions, and timeline. Must be handed in before the thesis topic is finally registered. \marginnote{\href{https://pages.schutera.com/1_Artifact/notes.pdf}{Artifact 1 Template}}
    \item \textbf{Midterm Report} (mid June) -- Thesis overview and proof of concept. Status against plan, overview figure of methods or architecture, prototype, and early challenges. \marginnote{\href{https://pages.schutera.com/2_Artifact/notes.pdf}{Artifact 2 Template}}
    \item \textbf{Submission Ready} (mid July) -- Guidelines and checklists. Structure, formatting, reproducibility, ethics, and scholarly writing verified before final submission. \marginnote{\href{https://pages.schutera.com/3_Artifact/notes.pdf}{Artifact 3 Template}}
\end{enumerate}


\section{The Rules that Bind}
\newthought{Two DHBW statutes govern your thesis} and are read together: the \emph{StuPrO Wirtschaft} sets the thesis specifics, the general \emph{DHBW StuPrO} the procedure. The essentials, each quoted to its paragraph:
\marginnote{\textbf{Sources} (cohorts from 1 October 2024):
\href{https://www.ravensburg.dhbw.de/fileadmin/user_upload/Dokumente/Amtliche_Bekanntmachungen/2025/27_2025_Bekanntmachung_DHBW_StuPrO_Wirtschaft_inkl._2._AES.pdf}{StuPrO Wirtschaft, Nr.~27/2025} and
\href{https://www.dhbw.de/fileadmin/user_upload/Dokumente/Amtliche_Bekanntmachungen/2025/21_2025_Bekanntmachung_DHBW_StuPrO_inkl_2._AES_01.pdf}{DHBW StuPrO, Nr.~21/2025}. The Wirtschaft statute prevails on conflict.}

\begin{itemize}\setlength\itemsep{0.2em}
    \item \textbf{Duration:} twelve weeks, in the fifth or sixth practice phase; start and end set by the Studiengangsleitung (StuPrO Wirtschaft \S5(2)).
    \item \textbf{Length:} 40 to 60 pages; figures, tables, directories, and appendices excluded; deviation needs supervisor consent (StuPrO Wirtschaft \S5(2)).
    \item \textbf{Credit:} 12 ECTS, a workload of 360 hours (StuPrO Wirtschaft \S5(2)).
    \item \textbf{Weight:} 20 percent of the final Bachelor grade (DHBW StuPrO \S46(3)).
    \item \textbf{Examiners:} one appointed supervisor also grades; a second examiner only if the first grades below 4.0, then the mean of both (StuPrO Wirtschaft \S4(2), \S4(4)).
    \item \textbf{Passing:} grade 4.0 (``ausreichend'') or better (StuPrO Wirtschaft \S4(3)).
    \item \textbf{Company role:} a qualified company person accompanies the work but cannot supervise or grade (StuPrO Wirtschaft \S5(2)).
    \item \textbf{Degree:} B.Sc. for Data Science und K\"unstliche Intelligenz, Wirtschaftsinformatik, and other named programs; otherwise B.A. (StuPrO Wirtschaft \S9).
    \item \textbf{Language:} German or English, per module description (DHBW StuPrO \S8).
\end{itemize}


\section{Eligibility and Initial Planning}
\newthought{Clarify internally within your company} whether a thesis topic is desired and feasible during the upcoming 12-week practice phase.
\marginnote{\textbf{Tip:} Discuss early with your manager.}

\section{Topic Coordination with Company and DHBW}
\newthought{Define a preliminary topic together with your company}. The topic should address a real business problem and be feasible within 12 weeks.
\marginnote{\textbf{Note:} Ensure alignment with DHBW professors.}

Alternatively, a topic may originate from DHBW projects. In this case, ensure approval by your company, as the thesis is written during the practice phase.

Identify and approach suitable DHBW professors whose expertise aligns with your topic. Apply for supervision early.
\marginnote{\textbf{Supervisor note (self):} Push the company to issue a Letter of Intent to the student, signaling interest and intent to continue collaboration beyond the thesis. Non-binding, but a strong signal of commitment. Raise it during topic coordination, not after submission.}

\section{Registration of the Thesis}
\newthought{Register your thesis via the DHBW Moodle system}. 
\marginnote{\textbf{Timeline:} Submit title 1 month before start.}

The provisional title must be submitted no earlier than one month before the official start date and no later than the beginning of the 12-week writing period.

\newthought{Formally}, you propose the topic together with your Dual Partner and register it in text form by the deadline the Studiengangsleitung sets; the approval reaches you before the processing time starts (DHBW StuPrO \S26(1)).

\newthought{Before the topic is finally registered}, hand in the research proposal (Artifact 1). It anchors the topic in literature and a clear problem statement, and gives your DHBW supervisor the basis to approve the final registration.
\marginnote{\textbf{Scientific Publication:} If co-authoring 
a scientific paper from the thesis matters to you, 
ask your company and industrial supervisor about options and how open they are if the work is exceptionally good.
A conference or journal submission touches data, methods, 
and IP they may want to gate, and the writing itself 
extends past the submission deadline. It adds effort on all three sides, 
so only push for it if it is something you want for yourself.}
The official start date marks the beginning of your writing period.

\section{Supervision and Writing Phase}
\newthought{During the 12-week writing period}, work closely with:
\begin{itemize}
    \item A company supervisor (technical support, no grading authority)
    \item A DHBW academic supervisor (responsible for evaluation and academic guidance)
\end{itemize}
\marginnote{\textbf{Tip:} Schedule regular meetings with your industrial supervisor.}

Ensure regular alignment meetings and early clarification of expectations regarding structure, methodology, and scope.

If sensitive company data is used, clarify the need for a confidentiality notice (Sperrvermerk) with your company supervisor.


\section{Communication Rhythm with me (Academic Side)}
\newthought{Throughout the 12-week writing period}, the rhythm is light and mostly yours to drive. Three artifacts mark your progress, and three meetings anchor our exchange. Count on the meetings; do not count on an email at every step.
\marginnote{\textbf{Flow:} Artifact 1 $\to$ Kickoff $\to$ Artifact 2 $\to$ Interim $\to$ Artifact 3 $\to$ Defense.}

\begin{enumerate}
    \item \textbf{Artifact 1, Research Proposal:} the first milestone you deliver, on the template: motivation, gap, research question, methodology sketch, and timeline.
    \item \textbf{Kickoff meeting:} align topic, scope, and milestones; bring a 3 to 5 sentence pitch and one sharp research question.
    \item \textbf{Artifact 2, Midterm Report:} an honest status update at mid-term: plan vs. reality, prototype, first results, and open risks.
    \item \textbf{Interim meeting:} reconcile status against the plan, show the prototype in five minutes, and adjust focus for the second half.
    \item \textbf{Artifact 3, Submission Ready:} the final stretch against the submission checklists: formatting, citations, reproducibility, Sperrvermerk, and language.
    \item \textbf{Defense:} present and defend the finished work. By then, the hard work is done.
\end{enumerate}

\newthought{You drive the contact.} Reach out when you are stuck or need a decision. I am responsive but light-touch: I will not necessarily flag each artifact or send reminders, and booking the three meetings is on you.
\marginnote{A short feedback survey stays open throughout. Two minutes from your side, whenever it helps, lets me tailor supervision to what you actually need.}

\section{Submission}
\newthought{Submit the thesis digitally via the official examination Moodle}. 
\marginnote{\textbf{Warning:} Late submission = automatic failure.}

\newthought{You hand in} one printed and one electronic copy to your Studienakademie, or electronic only where a suitable system is provided (DHBW StuPrO \S20(3)).

\newthought{Every thesis carries a declaration of independent authorship} in the statute's exact wording (DHBW StuPrO \S20(5)); its final sentence is dropped for electronic-only submission:
\begin{quote}
``Ich versichere hiermit, dass ich die vorliegende Arbeit selbstst\"andig verfasst und keine anderen als die angegebenen Quellen und Hilfsmittel verwendet habe und diese Arbeit bei keiner anderen Pr\"ufung mit gleichem oder vergleichbarem Inhalt vorgelegt habe und diese bislang nicht ver\"offentlicht wurde. Des Weiteren versichere ich, dass die eingereichte elektronische Fassung mit der gedruckten Ausfertigung \"ubereinstimmt.''
\end{quote}

\newthought{Document AI use in a tabular KI-Verzeichnis.} Where generative AI contributed to the work, the Fakult\"at Wirtschaft requires that a reader can tell which parts, text, figures, or tables, were AI-generated and to what extent. This is recorded in a tabular KI-Verzeichnis appended at the end of the thesis, listing the tools used and the type and extent of their use. Verbatim or paraphrased AI passages are additionally cited like any source, with their exports stored in the electronic appendix, and a signed Erkl\"arung zum Einsatz von KI follows directly after the Selbstst\"andigkeitserkl\"arung above. The exact form is set by your Studiengang (Richtlinien \S4.2.2), so confirm it with your supervisor. Like the declaration, the KI-Verzeichnis carries no page number and is not listed in the table of contents.
\marginnote{\textbf{Undocumented AI use} is treated as a T\"auschungsversuch, with sanctions up to ``nicht bestanden'' and de-registration.}

\newthought{If your company requires confidentiality}, add the statute's Sperrvermerk (DHBW StuPrO \S20(4)), which applies only when the Dual Partner demands it.
\marginnote{\textbf{Sperrvermerk:} ``Der Inhalt dieser Arbeit darf weder als Ganzes noch in Ausz\"ugen Personen au\ss{}erhalb des Pr\"ufungsprozesses und des Evaluationsverfahrens zug\"anglich gemacht werden, sofern keine anderslautende Genehmigung des Dualen Partners vorliegt.''}

\section{Defense (Colloquium)}
\newthought{Present and defend your thesis} in a colloquium. The minimum passing grade is 4.0.
\marginnote{\textbf{Format:} 20 min presentation, 10 min discussion.}


\section{When Things Go Sideways}
\newthought{Most theses run to plan.} For the rest, the general statute sets your options. Know them before you need them.

\newthought{A deadline extension} needs an application in text form, a temporary excused hindrance made credible with proof, and a statement from your Dual Partner. It is not granted for illness (DHBW StuPrO \S37).
\marginnote{\textbf{Extension vs. illness:} an extension covers external, temporary hindrances; illness runs through the make-up route instead.}

\newthought{If illness stops you}, report it to your Studienakademie without delay and prove it with a medical certificate on the DHBW forms. You then make up the examination rather than fail it (DHBW StuPrO \S34).

\newthought{Deception and plagiarism} count already as attempts. The DHBW may check your written work electronically, with sanctions up to exclusion and de-registration (DHBW StuPrO \S35).

\newthought{A failed thesis} can be repeated once, with a new topic to register within three months of the result. There is no second attempt (DHBW StuPrO \S39, \S40).

\newthought{To contest a grade}, request a review (\"Uberdenkungsverfahren) in text form, with concrete and comprehensible grounds pointing to a grading error (DHBW StuPrO \S45).


\clearpage
\section{Apply for Thesis Awards}
\newthought{Strong thesis work deserves recognition beyond the grade.} Several awards honor outstanding Bachelor's theses. A successful application sharpens your CV, opens doors to professional networks, and often carries prize money.
\marginnote{\textbf{Potential venues:}
Curated directory across German awards: \href{https://www.e-fellows.net/stipendien/preise-fuer-abschlussarbeiten}{e-fellows.net Preise f\"ur Abschlussarbeiten}.

\emph{You applied at a specific venue, let the future kids know!
Please add a link here via a pull request. This list lives from student contributions.}
}

\newthought{Timing matters.} Most awards have fixed annual deadlines, often months after submission. Prepare while the work is fresh. A typical application package bundles the thesis PDF, a one to two page abstract, your CV, and at least one supervisor endorsement.

\newthought{Check your NDA and Sperrvermerk first.} Many award procedures require submitting the full thesis PDF, and some publish the winning work. If your thesis is under a non-disclosure agreement with your industrial partner, confirm in writing that the agreement is permissive enough for the specific venue you are targeting. That includes scope, audience, and any publication by the award organizer. When in doubt, ask the partner's legal contact before you apply.

\newthought{Talk to me before you apply.} I can help identify the right venues for your topic, write the supervisor endorsement, and flag whether your industrial partner has a parallel internal route. If your work is strong, put it forward. Don't self-censor.
\marginnote{\textbf{Reach out early:} the endorsement letter alone usually needs one to two weeks of lead time.}

\newthought{Hall of fame.} Students who won an award on their thesis work appear in the margin. If that's you, please add yourself by pull request. Future cohorts see what's possible, and you get credit where it's due.
\marginnote{\textbf{Hall of fame} (name, year, thesis, company, venue):
\begin{itemize}\setlength\itemsep{0.2em}
    \item \emph{Your entry could live here.}
\end{itemize}
\emph{Won an award on your thesis? Please add yourself via a pull request.}
}



\iffalse
\chapter*{Roadmap for Master's Thesis}


\section{Eligibility to Start Writing}
\newthought{Contact the student office} to confirm you meet all academic requirements.

\marginnote{Contact: \\
\href{mailto:mail@dhbw.com}{mail@dhbw.com}}




\section{Topic Coordination with Supervisors}
\newthought{Before your thesis topic and supervision are approved, prepare a structured abstract}. 

\marginnote{Submission 1: \\
\href{https://pages.schutera.com/1_Artifact/notes.pdf}{Research Proposal Template}}
Send your abstract to your supervisor and schedule a meeting to refine the topic.



\section{Official Start of Thesis Writing}
\newthought{After topic approval,} register your thesis with the student office. The official cover page is then issued, marking the start of the writing period and submission deadline (MORE: 4 months, AMRD: 6 months).


\section{Thesis Submission}
Follow the official submission requirements from your student office and program coordinator.


\section{Defense}
Defend your thesis in a colloquium within six weeks of submission. The minimum passing grade is 4.0. Presentation: 20 minutes plus discussion.
\fi


% ----------------------------------------------------------------------



%\makestandardbackmatter{howtothesis_references}

\end{document}
